\ifdefined\gkver\else\def\gkver{student}\fi
\documentclass[\gkver]{gaokaozhenti}
\begin{document}

\chapter{2002年苏豫综合}
%% paperType: 文理综合物理部分

\begin{enumerate}
\item
%% number: 26
%% paperName: 2002年苏豫综合第 26 题 6 分
%% typeId: 1
%% type: 单选题
%% chapter: 气体性质
%% point: 热力学第一定律
%% method: 无
%% score: 6
%% degree: 750
%% duplicateId: 0
%% body:
一个带活塞的气缸内盛有一定量的气体。若此气体的温度随其内能的增大而升高，则\xzanswer{D}

\fourchoices[answer=D]
{将热量传给气体，其温度必升高}
{压缩气体，其温度必升高}
{压缩气体，同时气体向外界放热，其温度必不变}
{压缩气体，同时将热量传给气体，其温度必升高}

%% answer: D
%% memo:
\memoanswer{
本题原卷及材料中未提供详解，待补充。
}

\item
%% number: 27
%% paperName: 2002年苏豫综合第 27 题 6 分
%% typeId: 1
%% type: 单选题
%% chapter: 光学
%% point: 光的折射
%% method: 无
%% score: 6
%% degree: 750
%% duplicateId: 0
%% body:
已知一束单色光在水中的传播速度是真空中的 $\frac{3}{4}$，则\xzanswer{A}

\fourchoices[answer=A]
{这束光在水中传播时的波长为真空中的 $\frac{3}{4}$}
{这束光在水中传播时的频率为真空中的 $\frac{3}{4}$}
{对于这束光，水的折射率为 $\frac{3}{4}$}
{从水中射向水面的光线，一定可以进入空气中}

%% answer: A
%% memo:
\memoanswer{
本题原卷及材料中未提供详解，待补充。
}

\item
%% number: 28
%% paperName: 2002年苏豫综合第 28 题 6 分
%% typeId: 1
%% type: 单选题
%% chapter: 力物体的平衡
%% point: 物体系平衡
%% method: 整体与隔离
%% score: 6
%% degree: 650
%% duplicateId: 0
%% body:
如图所示，物体 a、b 和 c 叠放在水平桌面上，水平为 $F_{b}=5\UN$、$F_{c}=10\UN$ 分别作用于物体 b、c 上，a、b 和 c 仍保持静止。以 $f_{1}$、$f_{2}$、$f_{3}$ 分别表示 a 与 b、b 与 c、c 与桌面间的静摩擦力的大小，则\xzanswer{C}

\onepicture[w=5.5cm]{\includesvg{figs/28}}

\fourchoices[answer=C]
{$f_{1}=5\UN$，$f_{2}=0$，$f_{3}=5\UN$}
{$f_{1}=5\UN$，$f_{2}=5\UN$，$f_{3}=0$}
{$f_{1}=0$，$f_{2}=5\UN$，$f_{3}=5\UN$}
{$f_{1}=0$，$f_{2}=10\UN$，$f_{3}=5\UN$}

%% answer: C
%% memo:
\memoanswer{
本题原卷及材料中未提供详解，待补充。
}

\item
%% number: 29
%% paperName: 2002年苏豫综合第 29 题 6 分
%% typeId: 1
%% type: 单选题
%% chapter: 电场
%% point: 电势能电势差电场力做功
%% method: 无
%% score: 6
%% degree: 550
%% duplicateId: 0
%% body:
图中，a、b、c、d、e 五点在一直线上，b、c 两点间的距离等于 d、e 两点间的距离。在 a 点固定放置一个点电荷，带电量为 $+Q$，已知在 $+Q$ 的电场中 b、c 两点间的电势差为 $U$。将另一个点电荷 $+q$ 从 d 点移动到 e 点的过程中\xzanswer{D}

\onepicture[w=7.5cm]{figs/29.png}

\fourchoices[answer=D]
{电场力做功 $qU$}
{克服电场力做功 $qU$}
{电场力做功大于 $qU$}
{电场力做功小于 $qU$}

%% answer: D
%% memo:
\memoanswer{
本题原卷及材料中未提供详解，待补充。
}

\item
%% number: 30
%% paperName: 2002年苏豫综合第 30 题 6 分
%% typeId: 1
%% type: 单选题
%% chapter: 电磁感应
%% point: 电磁感应中的变加速
%% method: 无
%% score: 6
%% degree: 650
%% duplicateId: 0
%% body:
如图所示，在一均匀磁场中有一 U 形导线框 abcd，线框处于水平面内，磁场与线框平面垂直，$R$ 为一电阻，ef 为垂直于 ab 的一根导体杆，它可在 ab、cd 上无摩擦地滑动。杆 ef 及线框中导线的电阻都可不计。开始时，给 ef 一个向右的初速度，则\xzanswer{A}

\onepicture[w=5.5cm]{figs/30.png}

\fourchoices[answer=A]
{ef 将减速向右运动，但不是匀减速}
{ef 将匀减速向右运动，最后停止}
{ef 将匀速向右运动}
{ef 将往返运动}

%% answer: A
%% memo:
\memoanswer{
本题原卷及材料中未提供详解，待补充。
}

\item
%% number: 33
%% paperName: 2002年苏豫综合第 33 题 18 分
%% typeId: 3
%% type: 填空题
%% chapter: 原子和原子核
%% point: 质能方程
%% method: 无
%% score: 18
%% degree: 550
%% duplicateId: 0
%% body:
在其他能源中，核能具有能量密度大，地区适应性强的优势。在核电站中，核反应堆释放的核能被转化为电能。核反应堆的工作原理是利用中子轰击重核发生裂变反应，释放出大量核能。

\begin{enumerate}
	\item
	核反应方程式 \ce{^{235}_{92}U + n -> ^{141}_{56}Ba + ^{92}_{36}Kr + aX} 是反应堆中发生的许多核反应中的一种，n 为中子，X 为待求粒子，$a$ 为 X 的个数，则 X 为 \tkanswer[1cm]{\ce{n}}，$a=$\tkanswer[1cm]{3}。以 $m_{\mathrm{U}}$、$m_{\mathrm{Ba}}$、$m_{\mathrm{Kr}}$ 分别表示 \ce{^{235}_{92}U}、\ce{^{141}_{56}Ba}、\ce{^{92}_{36}Kr} 核的质量，$m_{\mathrm{n}}$、$m_{\mathrm{p}}$ 分别表示中子、质子的质量，$c$ 为光在真空中传播的速度，则在上述核反应过程中放出的核能 $\Delta E=$\tkanswer[3cm]{$(m_{\mathrm{U}}-m_{\mathrm{Ba}}-m_{\mathrm{Kr}}-2m_{\mathrm{n}})c^{2}$}。
	\item
	有一座发电能力为 $P=1.00\times10^{6}\UkW$ 的核电站，核能转化为电能的效率 $\eta=40\%$。假定反应堆中发生的裂变反应全是本题（1）中的核反应，已知每次核反应过程放出的核能 $\Delta E=2.78\times10^{-11}\UJ$，\ce{^{235}_{92}U} 核的质量 $m_{\mathrm{U}}=390\times10^{-27}\Ukg$。求每年（1 年 $=3.15\times10^{7}\Us$）消耗的 \ce{^{235}_{92}U} 的质量。
\end{enumerate}

%% answer:
\jdanswer{
\begin{enumerate}
	\item
	\ce{n}，3，$(m_{\mathrm{U}}-m_{\mathrm{Ba}}-m_{\mathrm{Kr}}-2m_{\mathrm{n}})c^{2}$
	\item
	$M=1.10\times10^{3}\Ukg$
\end{enumerate}
}

%% memo:
\memoanswer{
（1）由质量数守恒：$235+1=141+92+a$，解得 $a=3$；由电荷数守恒：$92=56+36$，故 X 为中子 \ce{n}。反应前总质量为 $m_{\mathrm{U}}+m_{\mathrm{n}}$，反应后总质量为 $m_{\mathrm{Ba}}+m_{\mathrm{Kr}}+3m_{\mathrm{n}}$，质量亏损 $\Delta m=m_{\mathrm{U}}+m_{\mathrm{n}}-m_{\mathrm{Ba}}-m_{\mathrm{Kr}}-3m_{\mathrm{n}}=m_{\mathrm{U}}-m_{\mathrm{Ba}}-m_{\mathrm{Kr}}-2m_{\mathrm{n}}$，放出的核能 $\Delta E=\Delta m c^{2}=(m_{\mathrm{U}}-m_{\mathrm{Ba}}-m_{\mathrm{Kr}}-2m_{\mathrm{n}})c^{2}$。

（2）反应堆每年提供的核能 $E_{\text{总}}=\frac{PT}{\eta}$（其中 $T$ 表示一年的时间）。

以 $M$ 表示每年消耗的 \ce{^{235}_{92}U} 的质量，得 $M:m_{\mathrm{U}}=E_{\text{总}}:\Delta E$。

解得 $M=\frac{m_{\mathrm{U}}PT}{\eta\Delta E}$。

代入数据得 $M=1.10\times10^{3}\Ukg$。
}

\end{enumerate}

\end{document}
